\documentclass[11pt]{article}

\usepackage{epsfig}
\usepackage{amsfonts}
\usepackage{amssymb}
\usepackage{amstext}
\usepackage{amsmath}
\usepackage{xspace}
\usepackage{theorem}
\usepackage{hyperref}
\usepackage{fullpage}
%\usepackage[]{algorithm2e}
\usepackage{xfrac}
\usepackage{mathtools}
\usepackage{algorithm}
\usepackage{algpseudocode}

\usepackage{listings}

\DeclarePairedDelimiter\ceil{\lceil}{\rceil}
\DeclarePairedDelimiter\floor{\lfloor}{\rfloor}

\usepackage{enumitem}                     

\usepackage{tikz}
\usepackage{tikz-qtree}
\usetikzlibrary{shapes}

\tikzstyle{code} = [black!90, draw=black!30, fill=black!5, very thick,
    rectangle, dashed, inner xsep=10pt, inner ysep=7pt]

\newenvironment{codebox}{
    \hspace{.05\textwidth}
        \begin{tikzpicture}
            \node[code] \bgroup
                \begin{minipage}{.65\textwidth}
                    \begin{alltt}}
                    {\end{alltt}
                \end{minipage}
            \egroup;
        \end{tikzpicture}
}

\newif\ifrubric
\rubrictrue
 
 
 
 \usepackage{titlesec}

\titleformat*{\section}{\bfseries}
\titleformat*{\subsection}{\bfseries}
\titleformat*{\subsubsection}{\bfseries}
\titleformat*{\paragraph}{\bfseries}
\titleformat*{\subparagraph}{\bfseries}


\newenvironment{proof}{{\bf Proof:  }}{\hfill\rule{2mm}{2mm}}
\newenvironment{proofof}[1]{{\bf Proof of #1:  }}{\hfill\rule{2mm}{2mm}}
\newenvironment{proofofnobox}[1]{{\bf#1:  }}{}
\newenvironment{example}{{\bf Example:  }}{\hfill\rule{2mm}{2mm}}

\newtheorem{fact}{Fact}
\newtheorem{lemma}[fact]{Lemma}
\newtheorem{theorem}[fact]{Theorem}
\newtheorem{definition}[fact]{Definition}
\newtheorem{corollary}[fact]{Corollary}
\newtheorem{proposition}[fact]{Proposition}
\newtheorem{claim}[fact]{Claim}
\newtheorem{exercise}[fact]{Exercise}

% math notation
\newcommand{\R}{\ensuremath{\mathbb R}}
\newcommand{\Z}{\ensuremath{\mathbb Z}}
\newcommand{\N}{\ensuremath{\mathbb N}}
\newcommand{\F}{\ensuremath{\mathcal F}}
\newcommand{\SymGrp}{\ensuremath{\mathfrak S}}

\newcommand{\size}[1]{\ensuremath{\left|#1\right|}}
%\newcommand{\ceil}[1]{\ensuremath{\left\lceil#1\right\rceil}}
%\newcommand{\floor}[1]{\ensuremath{\left\lfloor#1\right\rfloor}}
\newcommand{\poly}{\operatorname{poly}}
\newcommand{\polylog}{\operatorname{polylog}}

% anupam's abbreviations
\newcommand{\e}{\epsilon}
\newcommand{\half}{\ensuremath{\frac{1}{2}}}
\newcommand{\junk}[1]{}
\newcommand{\sse}{\subseteq}
\newcommand{\union}{\cup}
\newcommand{\meet}{\wedge}

\newcommand{\prob}[1]{\ensuremath{\text{{\bf Pr}$\left[#1\right]$}}}
\newcommand{\expct}[1]{\ensuremath{\text{{\bf E}$\left[#1\right]$}}}
\newcommand{\Event}{{\mathcal E}}
\newcommand{\E}{\ensuremath{\text{E}}}

\newcommand{\mnote}[1]{\normalmarginpar \marginpar{\tiny #1}}
\newcommand{\opt}[0]{{\text{OPT}}}

\setenumerate[0]{label=(\alph*)}


%%%%%%%%%%%%%%%%%%%%%%%%%%%%%%%%%%%%%%%%%%%%%%%%%%%%%%%%%%%%%%%%%%%%%%%%%%%
% Document begins here %%%%%%%%%%%%%%%%%%%%%%%%%%%%%%%%%%%%%%%%%%%%%%%%%%%%
%%%%%%%%%%%%%%%%%%%%%%%%%%%%%%%%%%%%%%%%%%%%%%%%%%%%%%%%%%%%%%%%%%%%%%%%%%%



\begin{document}

\noindent {\large {\bf 601.433/633 Introduction to Algorithms} \hfill {{\bf Fall 2026}}}\\ 
{{\bf Homework \#6}} \hfill {{\bf Due:} October 29, 2026, 11:59pm} \\
\rule[0.1in]{\textwidth}{0.4pt}

Remember: you may work in groups, but must write up your solution entirely on your own.  Solutions must by typeset (\LaTeX\ preferred but not required). 
Collaboration is limited to discussing the problems -- you may not look at, compare, reuse, etc.~any text from anyone else in the class. 
You may use the internet to look up formulas, definitions, etc., but may not simply look up the answers online.  

Please include proofs with all of your answers, unless stated otherwise.

\noindent \rule[0.1in]{\textwidth}{0.4pt}

\section{Computer Viruses (50 points)}
You're helping some security analysts monitor a computer network, tracking the spread of an online virus. There are $n$ computers in the system, 
labeled $C_1, C_2, \dots, C_n$. You are given a collection of \emph{trace data} indicating the times at which pairs of computers communicated. 
This data is a sequence of ordered triples $(C_i, C_j, t_k)$, where such a triple indicates that $C_i$ and $C_j$ exchanged bits at time $t_k$. 
There are $m$ triples total.

We'll assume that the triples are presented to you in sorted order of time, and for simplicity we'll assume that each pair of computers communicates 
at most once during the interval that you're observing.

The security analysts you're working with would like to able to answer questions of the following form: if the virus was inserted into computer $C_a$ 
at time $0$, what is the first time at which it could have have infected computer $C_b$? If it cannot be infected at all, 
then we say that this time is $\infty$.

The mechanics of infection are simple: if an infected computer $C_i$ communicates with an uninfected computer $C_j$ at time $T_k$ 
(i.e., if one of the triples $(C_i, C_j, t_k)$ or $(C_j, C_i, t_k)$ appears in the trace data), then computer $C_j$ becomes infected as well, 
starting at time $t_k$. Infection can thus spread from one machine to another across a \emph{sequence} of communications, 
provided that no step in this sequence involves a move backward in time. Thus, for example, if $C_i$ is infected by time $t_k$,
and the trace data contains triples $(C_i, C_j, t_k)$ and $(C_j, C_q, t_r)$ where $t_k \leq t_r$, then $C_q$ will become infected via $C_j$. 
Note that it is OK for $t_k$ to be equal to $t_r$; this would mean that $C_j$ had open connections to both $C_i$ and $C_q$ at the same time, 
and so a virus could move from $C_i$ to $C_q$ through $C_j$.

For example, suppose that $n = 4$ and the trace data consists of the triples

\begin{equation*}
    (C_1, C_2, 4), (C_2, C_4, 8), (C_3, C_4, 8), (C_1, C_4, 12)
\end{equation*}

and the virus was inserted into computer $C_1$ at time $0$.  Then $C_3$ would first be infected at time $8$ by a series of three steps: 
first $C_2$ becomes infected at time $4$, then $C_4$ from $C_2$ at time $8$, and then $C_3$ from $C_4$ at time $8$. 
So we should return $8$ if the query is $C_3$.

On the other hand, if the trace data were 

\begin{equation*}
    (C_2, C_3, 8), (C_1, C_4, 12), (C_1, C_2, 14)
\end{equation*}

and again the virus was inserted into $C_1$ at time $0$, then $C_3$ would \emph{not} become infected during the period of observation. 
So we should return $\infty$ if the query is $C_3$.

So your goal is to design an algorithm (as fast as possible) that answers questions of this type. The input is a collection of trace data 
(consisting of $m$ triples),  a virus injection point (a computer $C_a$), and another computer $C_b$. 
Your algorithm must return the first time at which $C_b$ could be infected, assuming the virus was injected into $C_a$ at time $0$.

As always, prove correctness and running time.

\bigbreak

Algorithm Description:

I use a weighted graph-based approach to solve this problem. We will use an undirected graph where the edges represent an interaction between $2$ 
computers, and the edge weight will represent the timestamp. We represent this graph using an adjacency list, where each vertex is 
mapped to those it's adjacent to with a tuple of the vertex and the time. The creation of the graph will clearly take $O(n + m)$ time, 
since it will first initialize each list for the $n$ computers (representing the vertices), and then iterate through the $m$ tuples 
and add each interaction to the graph representation.

From here, we must determine the shortest path -- that is, the fastest time it will take for $C_b$ to be infected given $C_a$ as the injection point. 
To do so, we will use a modifed version of Dijkstra's algorithm. Just like regular Dijksta, we maintain an array $d[n]$, representing the shortest 
distance from the source vertex, $C_a$. Initially, all values in the array are set to $\infty$. 
We maintain a minimum priority queue, initially containing $C_a$ with priority $0$. 

Normally, in Dijksta, during the relaxation, the weights are added together to form the complete path. For vertex $v$ communicating with vertex $u$, 
the shortest distance would be $d[v] = \min(d[v], d[u] + w(u, v))$. However, since the weight is the raw timestamp of communication, we can modify this: 
$d[v] = \min(d[v], \max(d[u], w(u, v)))$. This represents an infected computer $u$ spreading the infection to computer $v$ at a time after 
$u$ has become infected. 

The actual process of this algorithm would go as follows: Let $v$ represent the vertex with the smallest tentative extraction time (root of the queue). 
First, $v$ is extracted from the priority queue. Then, the algorithm examines each of its adjacent vertices, 
calculating the candidate infection time using the formula described above. If the candidate yields a smaller infection time, we update $d[v]$ 
and update the priority of $v$ in the queue. This repeats until either the queue is empty or $C_b$ is extracted. 
If $C_b$ is extracted, we return $d[b]$, since its infection time has been finalized. Otherwise, we return $\infty$, 
since no valid infection path exists from $C_a$ to $C_b$.

\bigbreak

Correctness Proof:

TODO

\bigbreak

Runtime Proof:

TODO

\section{Consistent Lifetimes (50 points)}

Suppose that you are an historian, and you are trying to figure out possible dates for when various historical figures may have lived. 
In particular, there are $n$ people $P_1, P_2, \dots, P_n$ who you are studying. Based on your research, 
you have discovered a collection of $m$ facts about when these people lived relative to each other. 
Each of these facts has one of the following two forms:

\begin{itemize}
\item For some $i$ and $j$, person $P_i$ died before person $P_j$ was born; or
\item For some $i$ and $j$, the lifetimes of $P_i$ and $P_j$ overlapped. 
\end{itemize}

Unfortunately, since the historical record is never fully trustworthy, it's possible that some of these facts are incorrect. 
Design an $O(n+m)$ time algorithm to at least check whether they are \emph{internally consistent}, i.e., 
whether there is a birth and a death date for each person so that all of the facts are true. 
As always, prove running time and correctness (i.e., prove that if there is a possible set of dates then your algorithm will return yes, 
and if no such dates are possible then your algorithm will return no).

\bigbreak

Algorithm Description:

We use an unweighted directed graph-based algorithm. For each person $P_i$ in the record, we create $2$ vertices in the graph -- 
$B_i$ (their birth) and $D_i$ (their death). This yields $2n$ vertices. We use directed edges to represent the chronology, 
so $B_i$ has an edge directed towards $D_i$ to represent the birth occurring before the death. 

We construct a graph based on the given $m$ facts. We initialize $2n$ adjacency lists to represent each known fact. 
Then, we iterate through the $m$ facts. We will create the edges as follows for each fact:

\begin{enumerate}[label=(\arabic*), noitemsep]
\item For some $i$ and $j$, person $P_i$ died before person $P_j$ was born: $D_i$ has an edge directed towards $B_j$.
\item For some $i$ and $j$, the lifetimes of $P_i$ and $P_j$ overlapped: $B_i$ has an edge directed towards $D_j$ and 
$B_j$ has an edge directed towards $B_i$.
\end{enumerate}

This construction process takes $O(2n + m)$, or $O(n + m)$ time. 

From here, we must detect directed cycles within the graph, since they show events that should take place after one another, which do not. 
We can do so using DFS with current path tracking. In this case, we keep track of $2$ arrays along the DFS: 
$visited[]$ and $path[]$. In this algorithm, for each unvisited vertex, we initialize a DFS traversal. 
Each vertex we visit will be marked \textbf{true} in both arrays. We then recurse through all neighbors. 
If we encounter a node already marked \textbf{true} in $path[]$, a cycle is detected and the algorithm returns \textbf{no}. 
After exploring all of the neighbors of any vertex $v$, we backtrack and reset $path[v]$ to \textbf{false}. 
If we complete the traversal with no cycle detection, the algorithm returns \textbf{yes}.

\bigbreak

Correctness Proof:

TODO

\bigbreak

Runtime Proof:

TODO

\end{document}
